% ============================================================================
% templates/exam/examen-oral/examen-oral.tex — plantilla: examen oral / sustentación
% ----------------------------------------------------------------------------
% Para: exámenes orales, defensas, sustentaciones y entrevistas académicas.
% Documento de uso del jurado: guion de la sesión, banco de preguntas guía
% y rúbrica de calificación individual.
% Uso: scripts/build.sh examen-oral.tex   (LuaLaTeX)
% ============================================================================
\documentclass{academic-exam}

\universidad{Universidad Nacional de San Cristóbal de Huamanga}
\facultad{Facultad de Ciencias Económicas, Administrativas y Contables}
\escuela{Escuela Profesional de Economía}
\curso{Proyectos de Inversión}
\codigocurso{EC-461}
\docente{Jurado: Mg. Nombre Apellido (presidente)}
\periodo{Semestre 2026-I}
\tipoevaluacion{Sustentación Oral}
\fechaevaluacion{21 de julio de 2026}
\duracion{25 minutos por estudiante}
\modalidad{Presencial · individual}

\begin{document}
\membrete

{\small
\renewcommand{\arraystretch}{1.7}
\begin{tabularx}{\textwidth}{@{}>{\bfseries\color{colorprimario}}l X
                              >{\bfseries\color{colorprimario}}l p{3.2cm}@{}}
  Estudiante: & \hrulefill & Código: & \hrulefill\\
  Proyecto sustentado: & \multicolumn{3}{@{}p{0.72\textwidth}@{}}{\hrulefill}\\
\end{tabularx}}
\par\vspace{0.7em}

\begin{instrucciones}[Protocolo de la sesión]
\begin{itemize}
  \item \textbf{Exposición} (10 minutos): el estudiante presenta el
        proyecto sin interrupciones. Al minuto 9 se avisa el cierre.
  \item \textbf{Ronda de preguntas} (10 minutos): cada miembro del jurado
        formula al menos una pregunta del banco guía o una repregunta.
  \item \textbf{Deliberación} (5 minutos): el jurado califica con la
        rúbrica adjunta, sin presencia del estudiante.
  \item La nota final es el promedio simple de los tres jurados,
        redondeado al medio punto.
\end{itemize}
\end{instrucciones}

% ===================================================== BANCO DE PREGUNTAS ==
\section*{\color{colorprimario}Preguntas guía del jurado}

\begin{pregunta}[][Dominio conceptual]
¿Qué criterios de decisión emplea su evaluación (VAN, TIR, ratio
beneficio-costo) y por qué el VAN es el criterio dominante cuando ambos
discrepan?
\end{pregunta}

\begin{pregunta}[][Supuestos]
Justifique la tasa de descuento utilizada. ¿Cómo cambiaría el resultado
si el costo de oportunidad del capital aumentara en dos puntos
porcentuales?
\end{pregunta}

\begin{pregunta}[][Análisis de riesgo]
Explique su análisis de sensibilidad: ¿cuáles son las dos variables más
críticas del proyecto y qué umbral las vuelve inviables?
\end{pregunta}

\begin{pregunta}[][Repregunta de profundización]
Si la demanda proyectada se redujera 20\,\%, ¿el proyecto seguiría siendo
socialmente rentable? ¿Qué evidencia de su estudio de mercado respalda su
respuesta?
\end{pregunta}

\begin{observacion}[Para el jurado]
Valore la \emph{calidad del razonamiento} antes que la memorización:
una respuesta que reconoce límites del propio trabajo y los discute con
rigor puntúa más alto que una defensa cerrada sin matices.
\end{observacion}

% ================================================================ RÚBRICA ==
\section*{\color{colorprimario}Rúbrica de calificación}

{\small
\renewcommand{\arraystretch}{1.45}
\arrayrulecolor{grislinea}
\begin{tabularx}{\textwidth}{@{}p{2.6cm} X X X c c@{}}
  \textbf{Criterio} & \textbf{Excelente (4)} & \textbf{Suficiente (2--3)}
    & \textbf{Insuficiente (0--1)} & \textbf{Peso} & \textbf{Nota}\\
  \hline\addlinespace[4pt]
  Dominio del tema &
    Explica con precisión conceptos, métodos y resultados; responde
    repreguntas sin apoyo. &
    Domina lo esencial; duda en repreguntas de profundización. &
    Errores conceptuales o incapacidad de sustentar resultados propios.
    & 35\,\% & \\
  \addlinespace[4pt]
  Argumentación &
    Razonamiento estructurado; distingue supuestos, evidencia y
    conclusiones. &
    Argumentos correctos pero incompletos o desordenados. &
    Afirmaciones sin sustento o contradictorias. & 25\,\% & \\
  \addlinespace[4pt]
  Claridad expositiva &
    Exposición ordenada, dentro del tiempo, con apoyo visual eficaz. &
    Exposición comprensible con problemas menores de ritmo o soporte. &
    Exposición confusa o fuera de tiempo. & 20\,\% & \\
  \addlinespace[4pt]
  Respuesta a preguntas &
    Responde de forma directa, honesta y fundamentada. &
    Responde parcialmente o con rodeos. &
    Evade o no responde. & 20\,\% & \\
  \addlinespace[6pt]\hline\addlinespace[6pt]
  \multicolumn{4}{r}{\textbf{Nota final (sobre 20):}} &
  \multicolumn{2}{l}{\hrulefill}\\
\end{tabularx}
\arrayrulecolor{black}}

\vspace{1.6em}
{\small
\begin{tabularx}{\textwidth}{@{}X X X@{}}
  \centering\rule{0.85\linewidth}{0.4pt}\par
  \centering\scshape Presidente del jurado &
  \centering\rule{0.85\linewidth}{0.4pt}\par
  \centering\scshape Primer miembro &
  \centering\rule{0.85\linewidth}{0.4pt}\par
  \centering\scshape Segundo miembro
  \tabularnewline
\end{tabularx}}

\findelexamen
\end{document}
